%%%

%%%   Самарский государственный технический университет

%%%   Кафедра прикладной математики и информатики

%%%

%%%   Преамбула для статьи в журнал "Вестник Самарского государственного

%%%   технического университета. Серия: «Физико-математические науки»"

%%%   Ориентирована под MikTEx 2.4 for Win32

%%%

%%%   Хотя сам журнал собирается под GNU/Linux на дистрибутиве TeXLive



\documentclass[11pt,a4paper,twoside]{article}



\usepackage[cp1251]{inputenc}

\usepackage[T2A]{fontenc}

\usepackage[english,russian]{babel}

\usepackage{samgtubib}

\usepackage[dvips]{graphicx}

\usepackage[multiple]{footmisc}



\usepackage{samgtu}

\renewcommand{\VestnikYear}{2009} % Год издания Вестника

\renewcommand{\VestnikNumber}{2\,(19)} % Номер Вестника

\renewcommand{\VestnikMonth}{Ноябрь} % Месяц выхода Вестника

\renewcommand{\VestnikData}{01.11.09} % Дата подписания Вестника в печать

\renewcommand{\VestnikUPL}{26,64} % Усл. печ. л.

\renewcommand{\VestnikUIL}{25,73} % Уч.-изд. л.

\renewcommand{\VestnikRegNumber}{479/09} % Регистрационный номер

\renewcommand{\VestnikOrder}{994} % Номер заказа






%
\usepackage{pst-barcode}
\usepackage{auto-pst-pdf}
%%
%\newcommand{\totop}[1]{\raisebox{6pt}[1pt][2pt]{#1}} 
%\newcommand{\tottop}[1]{\raisebox{12pt}[1pt][2pt]{#1}} 
%\newcommand{\totttop}[1]{\raisebox{18pt}[1pt][2pt]{#1}} 
%\newcommand{\tottttop}[1]{\raisebox{24pt}[1pt][2pt]{#1}} 
%\newcommand{\totttttop}[1]{\raisebox{30pt}[1pt][2pt]{#1}} 
%\newcommand{\tobot}[1]{\raisebox{-6pt}[1pt][2pt]{#1}} 
%\newcommand{\tobbot}[1]{\raisebox{-12pt}[1pt][2pt]{#1}} 
%\newcommand{\tobbbot}[1]{\raisebox{-18pt}[1pt][2pt]{#1}}


%\samgtupubl % для генерации pdf, отдаваемого в типографию СамГТУ
\pdfcrop % обрезка pdf для размещения на math-net.ru
%\draft % На корректуре
\filllastpage % Последняя страница не заполнена не полностью
%\briefcomm % Статья является кратким сообщением

\begin{document}



\hfill \begin{pspicture}(1in,1in)
\psbarcode{http://dx.doi.org/10.14498/vsgtu1465}{}{qrcode}
\end{pspicture}
\vspace{-1in}


\vsgtu{1465}

\FirstPage{799}
\LastPage{799}



\chead[\fancyplain{}{\small\sl Опечатка }]{\fancyplain{}{\small\sl Опечатка}}
  \rhead[]{\fancyplain{\Vestnik}{}}
  \lhead[\fancyplain{\Vestnik}{}]{}
%  \rhead[\fancyplain{\Vestnik}{}]{\fancyplain{\thepage}{\thepage}}
%  \lhead[\fancyplain{\thepage}{\thepage}]{\fancyplain{\Vestnik}{}}
  \cfoot[]{}
  \lfoot[\thepage]{}
  \rfoot[]{\thepage}
  \plainheadrulewidth=0.7pt
  \headrulewidth=0.4pt

  \pagestyle{fancyplain}

  \thispagestyle{plain}
  
  \addtocontents{toc}{\bigskip\par\hspace{-7mm}}
\addtocontents{etoc}{\bigskip\par\hspace{-7mm}}

\addcontentsline{toc}{section}{\it Дополнение к статье  Б.~О.~Волкова ``Даламбертианы Леви и их применение в квантовой теории''}

\addcontentsline{etoc}{section}{\it Supplement to the  article ``L\'{e}vy d'Alambertians and their application in the quantum theory'' by B.~O.~Volkov}


\renewcommand{\baselinestretch}{0.892}
\normalfont

{

$\;$
\bigskip

\bigskip

\bigskip

\begin{center}
{\Large \texttt{Дополнение к статье}}

\end{center}



\noindent
В о л к о в Б. О. Даламбертианы Леви и их применение в квантовой теории~// \textit{Вестн.
Сам. гос. техн. ун-та. Сер. Физ.-мат. науки}, 2015. T.~19, \No~2. С.~241–258
\crossref{10.14498/vsgtu1372}



\bigskip
\noindent
В статье Б.~О.~Волкова ``Даламбертианы Леви и их применение в квантовой теории'' место работы автора указано не полно. 
Второе место работы "---\\
{\small   \inst{2}\ruaffid{748}{Математический институт им.~В.~А.~Стеклова Российской академии наук,\\ 
Россия, 119991, Москва, ул. Губкина, 8}.}

\smallskip

\noindent
\textit{Сведения об авторе следует читать так}:

{\footnotesize \noindent
\fullauthor{\rupid{94935}{Борис Олегович Волков}}\regalia{к.ф.-м.н.; \mem{borisvolkov1986@gmail.com}} \position{доцент}\fdept{каф. ФН-12 <<Математическое моделирование>>\inst{1}} \lposition{старший научный сотрудник}\dept{отдел математической фи\-зики\inst{2}}
}




\bigskip


\EngVestnik

$\;$
\bigskip

\bigskip

\begin{center}
{\Large \texttt{Supplement to the article}}

\end{center}

\English

\noindent
V o l k o v B. O. L\'{e}vy d'Alambertians and their application in the quantum theory, \textit{Vestn. Samar. Gos. Tekhn. Univ., Ser. Fiz.-Mat. Nauki} [J. Samara State Tech. Univ., Ser. Phys. \& Math. Sci.],
2015, vol.~19, no.~2, pp.~241–258 \crossref{10.14498/vsgtu1372} (In Russian)



\bigskip


\noindent
In the article ``L\'{e}vy d'Alambertians and their application in the quantum theory'' by B.~O.~Volkov  the author affiliation are not fully specified. The second affiliation~is \\
{\small\inst{2}\enaffid{748}{Steklov Mathematical Institute of Russian Academy of Sciences, \\ 
Russia, 8  Gubkina st., Moscow, 119991, Russian Federation}.}



\smallskip

\noindent
\textit{Author Details  should read}:

{\footnotesize \noindent
\fullauthor{\enpid{94935}{Boris O. Volkov}}\regalia{Cand. Phys. \& Math. Sci.; \mem{borisvolkov1986@gmail.com}} \position{Assistant Professor}\fdept{Dept. of  Mathematical simulation\inst{1}} \lposition{Senior Researcher}\dept{Dept. of Mathematical Physics\inst{2}}
}


}

\Russian

\newpage

\end{document}
